% Problem record op_54aa8f0f61ecc5b1. This TeX file is derived from
% database/problems_json/op_54aa8f0f61ecc5b1.json by scripts/sync-tex.mjs; edit the JSON record
% and rerun the script rather than editing this file.
\section{Computability of ordinary quantum capacity}
\label{sec:op_54aa8f0f61ecc5b1}

\paragraph{Problem.}
Is the ordinary unassisted quantum capacity a computable function of a finite
description of a finite-dimensional quantum channel?  Let
$\mathcal N:\mathcal L(A)\to\mathcal L(B)$ be specified either by an exact
finite list of Kraus matrices with computable algebraic entries (including
rational entries) or by a finite quantum circuit over a fixed algebraic gate
set, allowing state preparation and partial trace.  For a complementary
channel $\mathcal N^c$, define coherent information and quantum capacity by
\begin{equation}
  I_c(\rho,\mathcal N)
  :=S(\mathcal N(\rho))-S(\mathcal N^c(\rho)),
  \qquad
  Q(\mathcal N)
  :=\sup_{n\geq1}\frac1n
       \max_{\rho_{A^{\otimes n}}}
       I_c(\rho_{A^{\otimes n}},\mathcal N^{\otimes n}),
  \label{eq:p72-quantum-capacity}
\end{equation}
where $S(\tau):=-\operatorname{Tr}(\tau\log_2\tau)$.  Writing
$\langle\mathcal N\rangle$ for either finite encoding above, the question is
whether there exists a Turing machine $T$ satisfying
\begin{equation}
  \forall\,\langle\mathcal N\rangle\ \forall k\in\mathbb N:
  \quad
  T(\langle\mathcal N\rangle,k)=q_{\mathcal N,k}\in\mathbb Q,
  \qquad
  |q_{\mathcal N,k}-Q(\mathcal N)|\leq2^{-k}.
  \label{eq:p72-computability}
\end{equation}
Determine whether the algorithm in Eq.~\eqref{eq:p72-computability} exists
for the capacity in Eq.~\eqref{eq:p72-quantum-capacity}, without imposing a
running-time bound.

\subsection*{Status}

Unsolved

\subsection*{Source}

Wilde explicitly asks whether quantum channel capacities are computable or
undecidable.  Bhattacharyya, Mehta, and Zhao formalize the complexity question
for succinct circuit descriptions and state that uncomputability of ordinary
quantum capacity remains unresolved
\sourcecite{ref:p72-wilde}{Wil17},
\sourcecite{ref:p72-bhattacharyya-mehta-zhao}{BMZ26}.

\subsection*{Progress}

\begin{itemize}
  \item Cubitt, Elkouss, Matthews, Ozols, P\'erez-Garc\'ia, and Strelchuk prove
  that no universal fixed tensor power detects all positive quantum
  capacities.  For every $m\in\mathbb N$, they construct a channel
  $\mathcal N_m$ such that
  \begin{equation}
    \max_{\rho_{A^{\otimes m}}}
      I_c(\rho_{A^{\otimes m}},\mathcal N_m^{\otimes m})=0,
    \qquad
    Q(\mathcal N_m)>0.
    \label{eq:p72-unbounded-uses}
  \end{equation}
  Equation~\eqref{eq:p72-unbounded-uses} rules out evaluating
  Eq.~\eqref{eq:p72-quantum-capacity} at one channel-independent blocklength,
  but it does not rule out a different terminating algorithm
  \sourcecite{ref:p72-cubitt-et-al}{CEM+15}.

  \item Bhattacharyya, Mehta, and Zhao prove that, for a channel supplied by a
  succinct quantum-circuit description, deciding whether
  $Q(\mathcal N)\geq3/4$ or $Q(\mathcal N)\leq1/4$, under that promise, is
  QMA-hard.  This is a lower bound on efficient computation, not a proof that
  the unrestricted Turing machine in Eq.~\eqref{eq:p72-computability} cannot
  exist \sourcecite{ref:p72-bhattacharyya-mehta-zhao}{BMZ26}.

  \item The same 2026 preprint proves an uncomputability result only for the
  restricted one-shot zero-error classical quantity
  $C^{(1)}_{0,\mathrm{PME}}$ of classical--quantum channels, where the shared
  state is maximally entangled and decoding is by a projective measurement.
  That theorem is not an ordinary-quantum-capacity result and therefore does
  not settle Eq.~\eqref{eq:p72-computability}.  The cited work is an
  unrefereed version~3 preprint; its authors report that version~2 contained
  an error in the undecidability section that version~3 corrects
  \sourcecite{ref:p72-bhattacharyya-mehta-zhao}{BMZ26}.
\end{itemize}

\subsection*{References}

\begin{enumerate}[label={},leftmargin=4.8em,labelsep=0.5em]
  \footnotesize
  \item[\textup{[Wil17]}]\label{ref:p72-wilde}
  M. M. Wilde, \emph{Quantum Information Theory}, 2nd ed., Cambridge
  University Press (2017), Sec.~26.6.
  \href{https://doi.org/10.1017/9781316809976}{doi:10.1017/9781316809976};
  \href{https://arxiv.org/abs/1106.1445}{arXiv:1106.1445}.

  \item[\textup{[CEM+15]}]\label{ref:p72-cubitt-et-al}
  T. Cubitt, D. Elkouss, W. Matthews, M. Ozols, D. P\'erez-Garc\'ia, and
  S. Strelchuk, ``Unbounded Number of Channel Uses May Be Required to Detect
  Quantum Capacity,'' \emph{Nature Communications} \textbf{6}, 6739 (2015).
  \href{https://doi.org/10.1038/ncomms7739}{doi:10.1038/ncomms7739};
  \href{https://arxiv.org/abs/1408.5115}{arXiv:1408.5115}.

  \item[\textup{[BMZ26]}]\label{ref:p72-bhattacharyya-mehta-zhao}
  A. Bhattacharyya, A. Mehta, and Y. Zhao,
  ``On the Undecidability of Quantum Channel Capacities,'' arXiv preprint
  (2026), version~3.
  \href{https://arxiv.org/abs/2601.22471v3}{arXiv:2601.22471v3}.
\end{enumerate}

\subsection*{Comment}

The precise gap is computability with no complexity bound, as formalized in
Eq.~\eqref{eq:p72-computability}.  QMA-hardness for succinct inputs, failure
of every universal fixed-block truncation, and uncomputability of a restricted
zero-error assisted capacity are all compatible with either answer to this
problem.

\subsection*{Field}

Quantum Communication; Quantum algorithm

\subsection*{Topic}

Quantum capacity; Computational complexity and computability

\subsection*{ID}

\texttt{op\_54aa8f0f61ecc5b1}
